\section{The contextual return fold}
\label{sec:return-fold}

All state numbers in this section are inherited labels. Let $O$ be the
pinned combined243 machine, with start state $95$, dispatcher $96$, and
halting state $254$. Its $243$ working states include $27$. A configuration
is a triple $(q,h,\tau)$, where $h\in\mathbb Z$ and
$\tau:\mathbb Z\to\{0,1\}$ is an arbitrary binary tape. Finite support is
not needed for the ordinary argument. Let $N$ be the machine obtained by
deleting state $27$ and applying precisely these four retained-row changes:
\[
\begin{array}{c|c|c}
\text{row}&O&N\\\hline
(15,0)&(0,R,21)&(0,L,32)\\
(18,0)&(0,L,27)&(0,L,15)\\
(47,0)&(0,L,27)&(0,L,15)\\
(65,0)&(0,L,27)&(0,L,15).
\end{array}
\]
Every other retained transition is unchanged. Thus $N$ has exactly $242$
working states and $484$ total binary transition rows. In particular,
$N$ is an ordinary two-symbol, one-tape, left/right machine; the saving
includes the blank-input initialization.

\paragraph{Inductive tape condition.}
For an $O$-configuration $c=(q,h,\tau)$, put
\[
 I(c)\quad\Longleftrightarrow\quad
 q\in\{15,18,47,65\}\ \Longrightarrow\ \tau(h-1)=0.
\]
The complete incoming-transition inventory of these states is
\[
\begin{array}{c|l}
\text{target}&\text{incoming source/read pairs}\\\hline
15&(11,0)\\
18&(12,0),(148,0)\\
47&(33,0),(144,0),(167,0)\\
65&(51,0),(81,0).
\end{array}
\]
Each listed transition writes zero and moves right. Consequently every
single step of $O$ that enters the displayed set establishes $I$; steps
ending elsewhere satisfy it vacuously. The initial state $95$ is outside
the set, so $I$ holds at the start on every initial tape and at every
initial head position. Induction on the number of steps proves $I$ at
all configurations reachable in $O$. This argument does not assume
anything about source-program registers or the outcome of the consistency
search.

\paragraph{The three-step replacement.}
Suppose $q=15$, $\tau(h)=0$, and $I(q,h,\tau)$. Thus the two cells at $h-1$
and $h$ are both zero. The old machine takes the single step
\[
 (15,h,\tau)\longrightarrow_O(21,h+1,\tau).
\]
The new machine takes the following three steps:
\[
 (15,h,\tau)\longrightarrow_N(32,h-1,\tau)
 \longrightarrow_N(45,h,\tau)
 \longrightarrow_N(21,h+1,\tau).
\]
The calculation uses the unchanged rows
$(32,0)\mapsto(0,R,45)$ and
$(45,0)\mapsto(0,R,21)$. Every write rewrites a zero as zero.
The endpoints therefore agree on the state, the head, and the entire
infinite tape. The zero immediately to the left supplies a return path
that recreates the old forward transition.

\paragraph{The three redirected entries.}
Let $q\in\{18,47,65\}$ and $\tau(h)=0$. The invariant again gives
$\tau(h-1)=0$. Using the old row
$(27,0)\mapsto(0,L,32)$, the two old steps are
\[
 (q,h,\tau)\longrightarrow_O(27,h-1,\tau)
 \longrightarrow_O(32,h-2,\tau).
\]
The new steps are
\[
 (q,h,\tau)\longrightarrow_N(15,h-1,\tau)
 \longrightarrow_N(32,h-2,\tau).
\]
Their complete endpoints are identical. No condition is imposed on
$\tau(h-2)$: state $32$ has not read that cell at the common endpoint.
The deleted state's read-one row is never used in these blocks.

\paragraph{All remaining rows.}
The complete incoming inventory of state $27$ consists of
$(18,0),(47,0),(65,0)$. Outside the four changed retained rows,
all $480$ rows in $N$ agree with the corresponding rows in $O$ and have
an $O$ target different from $27$. Consequently they give one old and
one new step to the same full configuration, with no deleted endpoint.
The two exceptional families above end in states $21$ or $32$, both
different from $27$ and from halt.

\paragraph{Global comparison.}
Take equal $O$ and $N$ configurations whose state is different from $27$
and satisfies the old-side invariant $I$. At a nonhalting checkpoint,
choose the block determined by the current state and scanned bit:
\[
\begin{array}{c|c|c|c}
\text{entry condition}&\text{old steps}&\text{new steps}
 &\text{common endpoint state}\\\hline
q=15,\ \tau(h)=0&1&3&21\\
q\in\{18,47,65\},\ \tau(h)=0&2&2&32\\
\text{all other nonhalting entries}&1&1&\text{the ordinary target}.
\end{array}
\]
The calculations prove equality of the next checkpoints. The invariant
holds at the old endpoint by its one-step inductiveness. The state is
again different from $27$. Every intermediate state in an exceptional
block is nonhalting, while an ordinary block reaches halt simultaneously
on its two sides.

Iteration gives a comparison from the shared state $95$. Both clocks
increase by at least one per block, so infinitely many comparison blocks
mean infinitely many actual transitions on both sides. Conversely, if
one side halts after finitely many transitions, it cannot pass through
infinitely many positive-duration blocks; the checkpoint comparison
therefore reaches the common halting state. This establishes
\[
 \operatorname{Halt}_O(95,h,\tau)
 \quad\Longleftrightarrow\quad
 \operatorname{Halt}_N(95,h,\tau)
\]
for every $h\in\mathbb Z$ and every binary tape $\tau$.
There is no consistency premise in this relative machine theorem.

The proof is a finite transition calculation followed by ordinary
induction. Under a conventional arithmetization of finite-support binary
configurations it is formalizable in PA: the block selectors and both
clock increments are bounded primitive-recursive operations, the local
cases have the fixed upper bound three, and the invariant is quantifier
free relative to the encoded tape lookup. This observation is an ordinary
formalizability argument, not a newly checked internal PA proof artifact.
The supplied executable checks the finite premises and an exact blank
prefix; it is not a proof-assistant kernel certificate.

\paragraph{Consistency calibration.}
Apply the comparison to the all-zero tape at head zero and compose it
with the inherited ordinary calibration of the exact combined243
machine. One obtains the ordinary theorem
\[
 \mathrm{ZFC}\ \vdash\
 \bigl(\neg\operatorname{Halt}(N)\ \longleftrightarrow\
 \operatorname{Con}(\mathrm{ZFC})\bigr).
\]
The new argument adds only the relative $243\to242$ machine comparison.
It does not change the logical frontend or the source theory, and it
does not close any previously open proof-assistant bridge in the inherited
source calibration. In particular, the claimed evidence level should not
be upgraded from an ordinary proof to a fully kernel-checked ZFC endpoint
merely because the local table reduction is elementary.

\paragraph{Reproduction and exact startup.}
The standalone driver reconstructs the table from the pinned old SHA256,
checks the complete incoming inventories, checks all unchanged rows, and
exhausts the sixteen relevant words on four cells. Every local block
reads only provided cells; the cells outside the word are arbitrary and
remain untouched. The local finite checks therefore verify the displayed
universal block calculations rather than testing a bounded collection
of source-register values.

As a separate execution check, the old blank prefix of $407\,967$ steps
and the new prefix of $407\,979$ steps reach exactly the same full
configuration: inherited state $96$, head $-5$, and the same set of $351$
one-cell positions. The old prefix visits $(15,0)$ six times. Each such
visit adds two steps under the proof, explaining the independently
observed twelve-step difference. This finite computation confirms the
binding and bootstrap; it is not evidence by itself for nonhalting.

\paragraph{Why the earlier source rotation enables this saving.}
This is not a direct fold of the earlier Routing245 table. In that
ancestor, row $(12,1)$ writes one and moves right into state $18$.
Such an entry leaves a one immediately left of state $18$, so the
displayed tape condition is not inductive for that ancestor. The QT source
rotation changes that routing. In combined243, the complete incoming
inventory of $18$ is precisely $(12,0)$ and $(148,0)$, both writing zero
and moving right. This makes state $27$'s read-one behavior unnecessary,
and the return through $32$ and $45$ provides a way to share its remaining
read-zero behavior with state $15$. The additional state saving thus
uses a tape invariant created by the source-level change, together with
a compensating change to the decoder. Older negative state-fold screens
on the unrotated table do not cover this construction.
